% v20: no zero low-D edge in a book low triangle; hence a nonzero D cycle.
% Depends on the canonical two-high form and the v19 finite-completion theorem.
\section{Book 低三角形の非零 \texorpdfstring{$D$}{D} 周期}
\label{bd20:section}

低辺の対向 $D$ 相殺は，node の向きとは別の符号の向きを定める．
零の $D$ を捨てて議論する前に，その場合を完全に排除する．

\begin{lemma}[低係数の零と負の隣接]
\label{bd20:local-zero}
正の整数二高分母標準形において，二つの低係数 $D_C,D_D$ が
同時に零となることはない．一方が負でもう一方が零となるのは，
canonical parameters
\begin{equation}\label{bd20:unique-zero}
 (r,s,z,S,T)=(4,4,5,2,2),\qquad (D_C,D_D)=(0,-2)
\end{equation}
の場合に限られる．この局所主張には分母の奇数性は不要である．
\end{lemma}

\begin{proof}
$\delta=T-S$ と置くと
\[
 D_C=\delta-\frac Sr s+\frac T{r+1}z,\qquad
 D_D=s-z-1.
\]
正の枝では $\delta\ge2$，$r\delta>S$ である．
$D_D<0$ なら整数性から $z\ge s$ であり，
$T/(r+1)>S/r$ によって $D_C>0$ となる．
$D_D=0$ のときは $s=z+1$ を直接代入して
\begin{equation}\label{bd20:positive-zero-identity}
 D_C=\frac{(r+z+1)(r\delta-S)}{r(r+1)}>0
\end{equation}
を得る．従って正の枝には負と零の組も，二つの零もない．

負の枝では $S=T=w$，$H=r-w$ とし，
\[
 c=\frac{sH}{r},\quad e=\frac{z(H+1)}{r+1},\qquad
 E=1-\frac{w^2}{H(H+1)}
\]
を用いる．既知の恒等式 $Ece=e-c+1$ と
$c\ge2,e\ge1,-1<E<1/2,E\ne0$ が成り立つ．
$E<0$ なら整数性から $c\ge e+2$ であり，
\[
 D_C=w\left(\frac e{H+1}-\frac cH\right)<0,
\]
\[
 s-z=(c-e)+w\left(\frac cH-\frac e{H+1}\right)>2.
\]
従って $D_D>1$ であり，零は生じない．

$E>0$ なら，定理~\ref{ml12:denominator-descent} の証明と同じく，
$j=e-c\ge1$，$c\le H$ となる．これらを再確認すると，
$e\le c-1$ は $Ece=e-c+1>0$ に反し，
$e=c$ は pair equation の $\gcd(c,e)=1$ と $c\ge2$ に反する．
また $E>0$ は整数性から $w\le H$ を与えるので
$E\ge1/(H+1)$，一方 $Ec=1-(c-1)/e<1$ だから $c\le H$ である．
従って
\[
 s-z=-j+\frac{w(c-Hj)}{H(H+1)}<0,\qquad D_D<0,
\]
\begin{equation}\label{bd20:negative-low-D}
 D_C=\frac{w(Hj-c)}{H(H+1)}\ge0.
\end{equation}
等号には $j=1,c=H$ が必要十分である．そのとき
$E=2/[H(H+1)]$ なので
\[
 w^2=H(H+1)-2.
\]
$H\ge3$ なら右辺は $H^2$ と $(H+1)^2$ の間に厳密にあるため，
$H=w=2$ に限られる．$c=2,e=3$ から $r=4,s=4,z=5$ を得る．
以上で零の全場合が尽くされた．
\end{proof}

\begin{lemma}[唯一の負・零 core は低三角形へ完成しない]
\label{bd20:zero-core-obstruction}
Book の仮定の下で，式~\eqref{bd20:unique-zero} の core を含む
低三角形は存在しない．
\end{lemma}

\begin{proof}
その core を $2$ とし，canonical order を $(0,1,3,4)$ とする．
行列と低係数は
\[
 P^2=
 \begin{pmatrix}
 0&-7&2&-4\\
 -14&0&-2&5\\
 4&-5&0&0\\
 -2&2&-1&0
 \end{pmatrix},\qquad
 (q_{20},q_{21},q_{23},q_{24})=(9,11,1,1),
\]
\[
 D^2_3=0,\qquad D^2_4=-2.
\]
低辺 $23$ の三つの node を $(n_0,n_1,n_4)$ とすると，
core $3$ の ray $2$ の行は
\[
 (-3-n_0,\ 6-n_1,\ 1-n_4),\qquad n_0+n_1+n_4=5
\]
に固定され，その $D$ は零である．
$n_4=1$ ならこの行は canonical $C$ で，高 ray の順序は $(1,0)$，
$2\le r\le6$ となる．負枝は $2\le w<r$ と $H$ の既知の範囲を，
正枝は $\delta=2,3$ と $r<S<r\delta$ を満たす．
$n_4=L+1$，$L=1,3,4$ なら canonical $D$ で，高順序は $(0,1)$ である．
$L=1$ では $S=T\in\{3,4,5,6\}$，
$L=3,4$ では可能な $(r,S,T)$ は
\[
 (2,3,5),\quad(3,4,6),\quad(2,3,6)
\]
に限られる．それ以外の行は二高標準形に合わない．

同様に低辺 $24$ の node を $(n_0,n_1,n_3)$ とすると，
core $4$ の ray $2$ の行は
\[
 (3-n_0,\ -1-n_1,\ 2-n_3),\qquad n_0+n_1+n_3=5
\]
で，その $D$ は二である．
canonical $C$ なら $n_3=2$，高順序 $(0,1)$，
$1\le r\le3$ である．canonical $D$ なら，可能なのは
高順序 $(1,0)$ の負枝 $w=2,3$ だけである．

以上の有限 reverse data に対し，指定した低 $D$ と高高 pair equation を連立する．
完全復元は定理~\ref{bk19:finite-completion} の二次式で実行できるが，
ここでは次の直接表示でも確認できる．
\[
 u=S+\frac Sr-1,\qquad v=T-\frac T{r+1}+1
\]
と置くと，高高 pair equation は
\[
 (uv-ST)sz+(u-S)s+(v-T)z+1=0.
\]
指定された低係数を $d$ とすると
\[
 z=
 \begin{cases}
 \displaystyle\frac{(r+1)S}{rT}s+\frac{(r+1)(d-\delta)}T,
       &\text{固定 ray が }C,\\[4pt]
 s-d-1,&\text{固定 ray が }D.
 \end{cases}
\]
この代入による二次式の整数根を取り，元の行列の整数性と奇数分母を検査する．
core $3$ の reverse data は六十八個，core $4$ では十五個であり，
全 scale の復元後に残る相手は次表の四個と一個だけである．
scale や高分母の上限は使っていない．

\begin{center}
\small
\begin{tabular}{cll}
相手 core & canonical $(r,s,z,S,T)$ & canonical order\\
\hline
$3$ & $(2,2,1,3,6)$ & $(0,1,4,2)$\\
$3$ & $(3,3,2,4,6)$ & $(0,1,4,2)$\\
$3$ & $(4,4,5,2,2)$ & $(1,0,2,4)$\\
$3$ & $(4,4,2,7,10)$ & $(1,0,2,4)$\\
\hline
$4$ & $(2,2,1,3,6)$ & $(0,1,2,3)$\\
\end{tabular}
\end{center}

残る低辺 $34$ を比べると，core $3$ の上表の順で次の値となる．
node の順序は punctures $(0,1,2)$ である．
\begin{center}
\small
\begin{tabular}{lcc}
core $3$ の canonical parameters & 最後の node 三つ & $D^3_4+D^4_3$\\
\hline
$(2,2,1,3,6)$ & $(2,-2,5)$ & $2$\\
$(3,3,2,4,6)$ & $(1,-1,5)$ & $1$\\
$(4,4,5,2,2)$ & $(2,-3,6)$ & $-2$\\
$(4,4,2,7,10)$ & $(-6,2,9)$ & $1$\\
\end{tabular}
\end{center}
いずれも node 非負性と $D$ 相殺に反する．
全二十一 node 行からの復元とこの表は
\nolinkurl{verify_book_low_D_cycle_v20.py} で独立に再計算し，
\nolinkurl{book_low_D_cycle_verification_v20.json} に保存した．
\end{proof}

\begin{theorem}[Book 低三角形の非零 $D$ 周期]
\label{bd20:nonzero-cycle}
正の整数局所 charts，奇数高分母，低辺の非負 node，
対向 $D$ 相殺を満たす book 低三角形が存在すると仮定する．
三本の低辺の全てで両端の $D$ は非零である．
各二価 core は一つの正の低 $D$ と一つの負の低 $D$ を持ち，
各低辺を負の端から正の端へ向けると有向三角形をなす．
\end{theorem}

\begin{proof}
二本以上の低辺が零の $D$ を持てば，その交点 core では二つの低 $D$ が
零になり，補題~\ref{bd20:local-zero} に反する．
零の辺がちょうど一本とする．その両端 core の残る低 $D$ は非零である．
両方が正なら第三 core の二つの低 $D$ が負になり，
補題~\ref{tl15:low-D} に反する．
従って少なくとも一方の端 core は負と零の組を持つ．
補題~\ref{bd20:local-zero} により式~\eqref{bd20:unique-zero} の模型に限られるが，
補題~\ref{bd20:zero-core-obstruction} がその完成を排除する．

これで全低 $D$ は非零である．三本の辺の相殺により，
六つの低端のうち正が三つ，負が三つとなる．
補題~\ref{tl15:low-D} により一つの core に負が二つ現れることはないので，
三つの cores はそれぞれ一つずつ負を持ち，残りは正である．
従って各頂点の入次数と出次数はともに一であり，有向三角形になる．
\end{proof}

この $D$ の向きは，高 ray $0$ への正 entry から定めた
補題~\ref{bk19:triangle} の向きとは別であり，
一致または反対のいずれも現時点では許している．
符号の周期性だけから book 全体の不存在は結論しない．
